%\documentclass[tikz]{standalone}
%\usepackage[american]{circuitikz}
%\begin{document}


%\begin{tikzpicture}[scale=1.5, transform shape,every node/.style={font=\footnotesize}]


\documentclass[tikz]{standalone}
\usepackage[american]{circuitikz}
\begin{document}
\begin{circuitikz}[american, scale=1, transform shape]
  %-------------------------
  % Nodo comune (Neutral) per la connessione a stella
  %-------------------------
  \coordinate (Neutral) at (0,0);
  
  %-------------------------
  % Induttanze che rappresentano i generatori (reti)
  %-------------------------
  % Fase A: da Neutral a punto A
  \draw (Neutral) to[L, l_={$L_A$}] (2,2) node[right] {A};
  % Fase B: da Neutral a punto B
  \draw (Neutral) to[L, l_={$L_B$}] (2,0) node[right] {B};
  % Fase C: da Neutral a punto C
  \draw (Neutral) to[L, l_={$L_C$}] (2,-2) node[right] {C};
  
  %-------------------------
  % Tiristori (SCR) disposti verticalmente verso il basso
  % Ogni tiristore parte dal rispettivo punto di fase e scende verso il bus DC+
  % Utilizzo del simbolo "scr" (silicon controlled rectifier) in circuitikz.
  %-------------------------
  % SCR per la fase A:
  \draw (2,2) -- (3,2); % collegamento orizzontale dalla fase al tiristore
  \draw (3,2) to[scr, rotate=90, l_={$SCR_A$}] (3,0);
  
  % SCR per la fase B:
  \draw (2,0) -- (3,0);
  \draw (3,0) to[scr, rotate=90, l_={$SCR_B$}] (3,-2);
  
  % SCR per la fase C:
  \draw (2,-2) -- (3,-2);
  \draw (3,-2) to[scr, rotate=90, l_={$SCR_C$}] (3,-4);
  
  %-------------------------
  % Formazione del bus DC positivo
  % I terminali inferiori (c

\end{circuitikz}
\end{document}





%\end{tikzpicture}


%\end{document}
